% Experiment report. Fixed sections: see experiments/TEMPLATE/report.tex and experiments/README.md.
% Build: npm run experiment:report -- 29-rl-econ-level
\documentclass{../polyexp}
\begin{document}
\input{generated}
\ExpHeader

\section{Question and hypothesis}
\textbf{The idea (David, 27--28 Nov 2026, issue \#166).} Keep the boost factor variable and put it in the state,
starting with economy. Instead of a one-turn boost (experiments 21--26), the economy weight becomes a lasting level
that the policy moves one step per turn. The level is in the state because it carries over: without it, the state
would not be Markov.

\textbf{Why economy, and why now.} Fixed schedules on the 2{,}000 check seeds ($600{,}000{,}000\dots$, before this
experiment; \texttt{rl-value eval} with turn-only policy files) show that r15 invests too much in economy late in the
game:

\begin{center}\small
\begin{tabular}{lrrr}
\toprule
economy factor (of r15's) & mean $d$ & better / worse & sign $p$ \\
\midrule
$\times0.5$ every turn          & $+309$ & 1{,}066 / 786   & $4\cdot10^{-11}$ \\
$\times1.5$ every turn          & $-313$ & 575 / 964       & 1.0 \\
$\times0.5$ in turns 0--9 only  & $+44$  & 552 / 545       & 0.43 \\
$\times2$ in turns 0--9 only    & $-129$ & 243 / 353       & 1.0 \\
$\times0.5$ from turn 15        & $+257$ & 1{,}024 / 616   & $3\cdot10^{-24}$ \\
$\times0.25$ from turn 15       & $+253$ & 1{,}066 / 638   & $1\cdot10^{-25}$ \\
$\times0.5$ from turn 25        & $+41$  & 258 / 127       & $1\cdot10^{-11}$ \\
\bottomrule
\end{tabular}
\end{center}

An economy investment pays back over the turns that are left, so late investments cannot pay back before turn 30. The
schedules were picked from 15 tried on the same seeds, so their best values are somewhat optimistic.

\textbf{Prediction} (before any run of this experiment):
\begin{enumerate}
  \item \emph{The learned schedule.} Without income in the state, the policy is one path of levels. It stays at
        $\times1$ or within one step of it through turn 9, and is at $\times0.5$ or lower from turn 20 on.
  \item \emph{It passes.} With the better state version, the chosen policy beats r15 on the evaluation seeds by
        $+150$ to $+350$ and passes the \#171 rule.
  \item \emph{As good as the hand schedule.} Against the hand schedule ($\times0.5$ from turn 16, written before the
        run), the policy is not significantly worse (mean $d \ge -50$).
  \item \emph{The model is roughly calibrated.} Its own predicted gain ($V(\text{start})$ minus the model's value of
        keeping $\times1$) is within 150 points of the measured gain. Each $Q$ value rests on thousands of games, and
        the successor of a transition is known up to income, so the maximisation bias of 25 and 26 is small.
  \item \emph{Income adds little.} The version with income is within $\pm 60$ points of the one without (paired),
        not significantly different.
\end{enumerate}
\textbf{Against the prediction:} if value iteration cuts economy early, or stays at $\times1$, while the fixed
schedules gain $+250$, the model is misled. The likely cause would be that the state lacks what an early investment
builds up, so later rewards are credited to the later levels.

\section{Mechanism}
\textbf{Action set \texttt{econ5}} (\texttt{rl\_value::strategic}): the base bot's build factor (\texttt{type.build},
move type 13) times 0.25, 0.5, 0.75, 1 or 1.5; at $\times1$ the weights are r15's, bit for bit (a test checks). The
level starts at $\times1$ and moves at most one step per learner turn; the action is the level played this turn. A
state without a policy entry keeps the level.

\textbf{State} (\texttt{level\_key}): (turn, level before the turn) and, as a second version, (turn, level, income
/ 3). Without income the next state is fixed by the action. With income, the next income bucket is random.

\textbf{Value iteration} (\texttt{rl-value fit-level}, new): certainty-equivalence model as in 21--26, $Q(s,a) =$
mean over observed transitions of $r + V(s')$, with $r$ the score gained during the turn and $\gamma = 1$. $V(s) =
\max_a Q(s,a)$ over the actions with at least $m$ samples; otherwise the mean over all of $s$'s transitions, and
the level is kept. The turn is in the state, so one backward sweep is exact. The same model also evaluates ``keep
$\times1$ all game'' (r15) for the predicted gain.

\textbf{Data} (\texttt{--driver walk}): a random walk over the levels from $\times1$, each turn one step down, stay
or one step up with probability $1/3$ each (fixed by seed and turn), held at the ends.

\textbf{Playing a level policy} keeps the level between turns (\texttt{StrategicPolicy::picker},
\texttt{Driver::picker}); \texttt{rl-gate}'s \texttt{rl-policy:} learner uses the same picker. On 300 evaluation seeds
($10^9\dots$), \texttt{rl-value eval} and \texttt{rl-gate} give identical scores for a hand-written level policy
(\texttt{schedule-hand-x0.5-from-t16.tsv}: $\times0.75$ at turn 15, $\times0.5$ from turn 16): candidate mean
7{,}346.2, base 7{,}047.7, $d = 298.483$ in both. That check used 300 of the evaluation seeds; the hand schedule was
written from the check-seed table above, not tuned on them.

\section{Setup}
As experiments 25 and 26: learner base \texttt{greedy-full-probe-perf-r15}; opponents r15, probe-perf and the anchor
probe-perf (\texttt{rl-gate} defaults after PR \#204), shuffled per seed. Perfection, 16$\times$16 Continents,
4 players, Crazy, learner Xin-xi in the human seat. Seeds: training $1\dots300{,}000$ (random walk), action check
$600{,}000{,}000\dots$ (2{,}000; each level every turn), validation $700{,}000{,}000\dots$ (2{,}000), evaluation
$10^9\dots$ (10{,}000).

Commands (from \texttt{engine/}):
\begin{verbatim}
B=./target/release/rl-value; R=target/rl-value; E=../experiments/29-rl-econ-level
$B check --actions econ5 --games 2000 --seed-start 600000000 --out $E/check.json
$B collect --actions econ5 --driver walk --games 300000 --seed-start 1 --out $R/econ5-walk.tsv
$B fit-level --train $R/econ5-walk.tsv --income-bucket 0 --min-count 20,100,500 --out $R/level
$B fit-level --train $R/econ5-walk.tsv --income-bucket 3 --min-count 20,100,500 --out $R/level
$B eval --actions econ5 --policy $R/level/policy-bB-mM.tsv --n 2000 --seed-start 700000000 --out $E/val-bB-mM.json
rl-gate --candidate rl-policy:experiments/29-rl-econ-level/policy-bB-mM.tsv --base greedy-full-probe-perf-r15
rl-gate --candidate rl-policy:experiments/29-rl-econ-level/policy-bB-mM.tsv \
  --base rl-policy:experiments/29-rl-econ-level/schedule-hand-x0.5-from-t16.tsv --no-info15
\end{verbatim}

\section{Results}
\ExpResults

\todo{claude}{After the runs.}

\section{Behaviour}
\todo{claude}{The learned schedule (mean factor by turn) of both versions.}
\todo{jakob}{What you saw in the live games, if you watched any.}

\section{Interpretation}
\todo{claude}{Why the result came out this way, and whether the prediction held.}
\todo{jakob}{Your reading: agree, disagree, what is missing?}

\section{Threats to validity}
\todo{claude}{Two or three concrete threats.}

\section{Verdict and next step}
\todo{claude}{Kept or not, and the next experiment this suggests.}

\section{Notes for the paper}
\todo{claude}{Draft: the one-sentence claim and which figure or table would carry it.}
\todo{jakob}{Is this worth a place in the paper? Where?}

\end{document}
